\subsection{Keep each assessment attached to its historical terms}
\label{sec:nested-term-invariants}

A claim can remain outstanding while its due date changes. Checking its
current balance and current due date is therefore insufficient to explain an
earlier overdue assessment. In the retained synthetic example, a claim issued
on 2 January 2027 is initially due on 31 January. An imported assessment dated
28 February records 20 remaining and 20 overdue. Prospective extensions become
effective on 10 March and 10 April, moving the due date first to 30 April and
then to 31 May. Applying the current 31 May due date to the February assessment
would incorrectly turn its overdue amount into zero. The historical assessment
must remain attached to the initial term.

The term history is stored inside the claim. Each term has its own local
identifier, effective date, due date, predecessor and event identity. Each
assessment has a date, referenced term, remaining amount, overdue amount,
origin and event identity. Local identifiers alone do not identify these
records: another claim may use the same identifier. Nor does insertion order
establish chronology. A declaration must name the parent claim and select a
term using its effective date before comparing the financial amounts.

\paragraph{Expose the supplied rows without choosing a history.}
\code{NestedMapProjection/1} binds a declared parent map and nested child map
to typed scalar rows. Bindings may read the parent key, child key, a parent
field or a child field. The accompanying lineage retains both raw keys and
escaped source paths. Each parent also has a group record with its child-map
path, first row and row count. An empty child map still produces a group;
an empty parent map produces no groups. This distinction lets a rule require
terms for every supplied claim even when one claim has no terms.

Projection preserves supplied insertion order and makes no chronological
selection. Declarations are checked against the memory definition even when
there are no rows. An explicit population declaration identifies the exact
book and both map paths; partial or unknown coverage remains unresolved.
Complete supplied internal maps are sufficient for these checks, but do not
establish that an issuer's disclosures reveal all contracts or past events.

\paragraph{Select within one parent and one declared boundary.}
\code{GroupedRowQuery/1} provides count, typed match and greatest-row queries.
The query names the parent key, escaped parent path and child-map path; all
must agree. Count includes all rows in that group. Match compares one declared
field with one value of exactly the same type. Greatest orders dates or exact
rational values, optionally below an explicitly inclusive or exclusive upper
bound. Units, scale and context must agree before a rational comparison.

The February assessment selects the greatest effective date no later than
28 February, which is 2 January. A predecessor query instead uses a strict
upper bound: the predecessor of the term effective on 10 April is the term
effective on 10 March. A query without an upper bound identifies the latest
term. None of these operations sorts or replaces the stored source order.
Candidate and winner indices retain that order, and a unique result includes
the whole selected row and its source lineage.

Equal greatest dates produce an ambiguous result. Duplicate matches do the
same; insertion order is never a tie breaker. A group absent from a complete
supplied population is invalid. A present group with no qualifying row returns
explicit \code{none} when the declaration permits it, or invalid when the row
is required. An explicit unknown input remains unresolved; omitting a required
input is a malformed declaration. Every evaluation recomputes its projection
from the bound book, schemas, types and coverage. A caller cannot substitute
previously computed rows after changing the source.

\paragraph{Use absence without treating uncertainty as absence.}
\code{QueryGuardedInvariant/1} binds query results into the preserved guarded
requirement evaluator. Count becomes an exact dimensionless rational. A unique
row supplies \code{has\_row=true} and known selected fields. Allowed
\code{none} supplies \code{has\_row=false}, while its selected fields remain
unknown. An ambiguous or unresolved query supplies unknown values, including
an unknown presence indicator. Thus a tie cannot masquerade as an empty group.
The declaration explicitly chooses whether ambiguity is invalid or unresolved.

A false guard makes a predicate inactive; it does not prove that predicate.
The separate \code{required\_known} declaration can require a selected field
even when its predicate is inactive. Invalid queries fail structurally before
the predicates run, and unresolved queries keep the overall result unresolved
even under a false guard. Allowed absence is different: it can legitimately
activate an empty-group rule while leaving a selected-row rule inactive. The
result retains the original guard findings and distinguishes
\code{no\_active\_requirements} from satisfied requirements.

\paragraph{Declare the rule once for every supplied claim and assessment.}
Finite queries with handwritten claim identities would miss newly supplied
records. \code{NestedQueryInvariantTemplate/1} instead iterates over all groups
or all child rows of a declared complete projection. Typed captures read the
current parent or child identity, scalar fields and explicit book dates or
memory values. Every query's full group identity is bound to that parent.
Scalar substitutions can fill only a typed match value, typed bound or the
value of a typed constant already present in a guarded predicate. They cannot
replace operators, types, bindings or failure consequences. Duplicate or
unsafe substitutions refuse before iteration, including on empty populations.

Each instance retains its captured values, source paths, resulting declaration
and complete query and predicate results. The reviewed selection identifies
the template as a whole. The instance selection is explicitly derived by the
checked substitution procedure; it is not described as an independent review
of every resulting instance. The initial checked profile permits at most
64 instances and refuses overflow without a partially successful result.
This bound is an implementation choice for the finite probe, not a restriction
on the financial concepts that the product must eventually represent.

The term-row declaration uses a unique match at each effective date to reject
ties. It selects the greatest strictly earlier term and checks the stored
predecessor against that identity; when no predecessor exists, the explicit
initial sentinel must be \code{none}. It also checks issue, effective, due and
book dates. The assessment-row declaration checks uniqueness of its date,
selects the greatest applicable term, preserves historical overdue, requires
nonnegative assessed remaining amounts and distinguishes imported snapshots
from recorded events. Parent-group declarations check current terms, latest
assessments, assessed flags, summary amounts and empty-assessment sentinels.

Event membership uses the existing collection-reference evaluator through
\code{ProjectedStateCollections/1}. It binds complete term rows and the book's
seen-event sequence. An explicitly typed equality filter selects assessments
whose origin is \code{recorded\_event}; their event identities must be seen.
Imported snapshots have a separately declared unknown-event condition. A
filtered collection is complete only relative to its complete supplied source
and explicit predicate. The filter does not turn a partial source into a
complete collection, and its types are checked even when it selects no rows.

\paragraph{Join the semantic checks and preserve the admission boundary.}
\code{TermStateSemanticBundle/1} joins six reviewed components: common claim
conditions and the receivables total, the collection-clock bound, term history,
historical assessments, parent summaries and event membership. The clock reads
the actual exclusive fiscal endpoint rather than assuming a calendar year.
The common claim declaration retains nonnegative remaining amounts, sale-event
membership, issue dates and the exact receivables sum. Historical assessment
rules own overdue; the earlier scalar rule based on current due is not applied
to these amended-term histories.

All 22 retained states satisfy this semantic bundle. Fifty-one complete
comparisons preserve every subassessment, and 27 targeted semantic corruptions
are rejected by both the independently run native validator and the declared
checks. The candidate evaluations disable native financial callbacks. A
separate discriminator changes the native migration premise: native validation
rejects that premise while the book's declared semantic relations still hold.
This difference identifies the remaining boundary rather than an alternative
form of opening admission. Exact native structural schemas, role and cost
selection, and migration premises still need their own selected contract.
The bundle grants no execution or accounting admission and is not yet bound
into the opening-to-program execution path. Direct source authoring for these
nested profiles also remains to be supplied.

Complete declarations, original-state comparisons, structural controls and
the native-oracle outcomes are retained under
\begingroup\raggedright
\path{data/accounting-master/grouped-row-query-20261026/},
\path{data/accounting-master/query-guard-composition-20261026/},
\path{data/accounting-master/nested-query-template-20261026/},
\path{data/accounting-master/nested-obligation-completion-20261026/}
and \path{data/accounting-master/term-semantic-composition-20261026/}.
\par\endgroup
